\section{Introduction}
\noindent
Let $\K$ be a field. A Koszul algebra is a standard graded $\K$-algebra such that $\K$ has a linear minimal free resolution. The notion of Koszulness was first introduced by Priddy \cite{Priddy1970}. Given the infinite and unpredictable nature of a free resolution, it is quite difficult to verify Koszulness directly from its definition. One must instead use different methods. One method is to show that, for some presentation, the defining ideal of the algebra has a quadratic Gröbner basis. If this is the case, we call the algebra \emph{G-quadratic} \cite[Remark 2.6(4)]{conca2012koszulalgebrasregularity}. Another important method is to find a family of ideals within the algebra that forms a \emph{Koszul filtration} \cite{CTV2001}. Both G-quadraticity and the existence of a Koszul filtration are sufficient conditions for Koszulness, but neither is necessary \cite{CTV2001}.

One particularly elusive Koszul algebra is the cubic pinched Veronese algebra.
The pinched Veronese is defined as \[PV=\K [x^3,x^2y,xy^2,x^2z,xz^2,y^3,y^2z,yz^2,z^3]\subset \K[x,y,z]\] where $\K$ is a field. 
In 1993, Bernd Sturmfels asked Irena Peeva whether the pinched Veronese is Koszul and more than a decade later, Caviglia answered this question affirmatively \cite{Caviglia2009}. A more theoretical proof was subsequently given by Caviglia and Conca \cite{CC2013}. Both proofs had to go outside the more standard methods described above.

It is still unknown whether the pinched Veronese is G-quadratic, although one might conjecture the answer to be no. Using a Gröbner fan as first introduced by Mora and Robbiano \cite{mora_grobner_1988}, Conca, De Negri and Rossi showed that for toric coordinates it has no quadratic Gröbner basis \cite{conca2012koszulalgebrasregularity}. It is also unknown if it has a quadratic Gröbner basis after a lifting by some linear form, i.e. if it is \emph{LG-quadratic}. 

It was previously shown by this author along with Nicklasson, Thompson and Westerbäck that the pinched Veronese does not have a Koszul filtration by monomial ideals \cite{berghofer2026con}. It was previously unknown, however, whether it has a Koszul filtration in general. In this paper, we will add another layer to the mystery of the pinched Veronese as we show that despite it being Koszul, it does not have a Koszul filtration regardless of $\K$.



\begin{theorem*}\label{bigone}The pinched Veronese does not have a Koszul filtration for any~field~$\K$.
\end{theorem*}